% Practice list D --- topics of Lista 7 (Benigno 2015, AS-AD New Keynesian).
% \solucoesfalse -> blank list ; \solucoestrue -> answer key + rubric.
% Compile: pdflatex (2x); check: pdffonts (all Type 1). Numbers: practice-lists-check.py
\documentclass[11pt,a4paper]{article}
% Shared preamble of the practice lists (practice-L*.tex), adapted from lista-exam.tex.
% Each list sets \solucoesfalse / \solucoestrue and \listatema right after % Shared preamble of the practice lists (practice-L*.tex), adapted from lista-exam.tex.
% Each list sets \solucoesfalse / \solucoestrue and \listatema right after % Shared preamble of the practice lists (practice-L*.tex), adapted from lista-exam.tex.
% Each list sets \solucoesfalse / \solucoestrue and \listatema right after \input{practice-preamble}.
\usepackage[utf8]{inputenc}
\usepackage[T1]{fontenc}
\usepackage{lmodern}   % Type 1 vector fonts (replaces PK bitmaps)
\usepackage{cmap}      % ToUnicode CMap -> copyable text and accents
\usepackage[english]{babel}
\usepackage{amsmath,amssymb,amsfonts}
\usepackage[a4paper,margin=2.2cm]{geometry}
\usepackage{enumitem}
\usepackage{fancyhdr}
\usepackage{comment}
\usepackage{xcolor}
\usepackage{microtype}
\usepackage{booktabs}
\usepackage{tikz}

\newif\ifsolucoes
\newcommand{\disciplina}{Macroeconomics I --- MPE Insper}
\newcommand{\espacoresol}{4.5cm}   % height of the answer space in the blank list
% The sol environment must be decided after the list sets its toggle.
\AtBeginDocument{\ifsolucoes\includecomment{sol}\else\excludecomment{sol}\fi}

\pagestyle{fancy}\fancyhf{}
\lhead{\footnotesize\disciplina}
\rhead{\footnotesize Practice list\ifsolucoes{} --- Solutions and rubric\fi}
\cfoot{\footnotesize\thepage}
\renewcommand{\headrulewidth}{0.3pt}
\setlength{\headheight}{14pt}
\setlist{topsep=2pt,itemsep=2pt}
\newcounter{questao}

\newcommand{\bloco}[3]{\bigskip\noindent
  {\large\bfseries\color{blue!55!black}Block #1 --- #2}\hfill\textit{(#3)}\par
  \rule{\linewidth}{0.3pt}\smallskip}
\newcommand{\questao}[1]{\refstepcounter{questao}\medskip\noindent
  \textbf{Question \thequestao.}~#1\par}
\newenvironment{alternativas}{\begin{enumerate}[label=\Alph*.,leftmargin=2.2em,topsep=1pt]}{\end{enumerate}}
\newcommand{\gabarito}[1]{\ifsolucoes\par\noindent\textbf{Answer:}~#1\fi}
\newcommand{\resolespaco}{\ifsolucoes\else\par\textit{Solution:}\vspace{\espacoresol}\par\fi}
\newcommand{\passo}[2]{\par\smallskip\noindent\textbf{Step --- #1.}~#2}
\newcommand{\intuicao}[1]{\par\smallskip\noindent\textit{Economic intuition.}~#1}
\newcommand{\refbib}[1]{\par\smallskip\noindent{\footnotesize\color{black!60}\textbf{Ref:} #1}}
\newcommand{\rubrica}[1]{\par\smallskip\noindent\textbf{Grading rubric.}\begin{itemize}[leftmargin=1.4em,topsep=1pt]#1\end{itemize}}

\newcommand{\cabecalho}{%
\begin{center}
  {\Large\bfseries \disciplina}\\[2pt]
  {\large Practice list: \listatema}\ifsolucoes\\[2pt]\textit{Solutions and grading rubric}\fi
\end{center}
\vspace{2pt}
{\footnotesize\noindent\textbf{Instructions.} Multiple choice: mark the single correct option (no
penalty for wrong answers). \emph{True/False}: every 2 wrong items cancel 1 right one
(score $=\max\{0,\text{right}-\text{wrong}/2\}\times$ item value); blank items do not count.
Discursive questions: show every step. Numbers are illustrative and chosen to give clean answers.\par}}
.
\usepackage[utf8]{inputenc}
\usepackage[T1]{fontenc}
\usepackage{lmodern}   % Type 1 vector fonts (replaces PK bitmaps)
\usepackage{cmap}      % ToUnicode CMap -> copyable text and accents
\usepackage[english]{babel}
\usepackage{amsmath,amssymb,amsfonts}
\usepackage[a4paper,margin=2.2cm]{geometry}
\usepackage{enumitem}
\usepackage{fancyhdr}
\usepackage{comment}
\usepackage{xcolor}
\usepackage{microtype}
\usepackage{booktabs}
\usepackage{tikz}

\newif\ifsolucoes
\newcommand{\disciplina}{Macroeconomics I --- MPE Insper}
\newcommand{\espacoresol}{4.5cm}   % height of the answer space in the blank list
% The sol environment must be decided after the list sets its toggle.
\AtBeginDocument{\ifsolucoes\includecomment{sol}\else\excludecomment{sol}\fi}

\pagestyle{fancy}\fancyhf{}
\lhead{\footnotesize\disciplina}
\rhead{\footnotesize Practice list\ifsolucoes{} --- Solutions and rubric\fi}
\cfoot{\footnotesize\thepage}
\renewcommand{\headrulewidth}{0.3pt}
\setlength{\headheight}{14pt}
\setlist{topsep=2pt,itemsep=2pt}
\newcounter{questao}

\newcommand{\bloco}[3]{\bigskip\noindent
  {\large\bfseries\color{blue!55!black}Block #1 --- #2}\hfill\textit{(#3)}\par
  \rule{\linewidth}{0.3pt}\smallskip}
\newcommand{\questao}[1]{\refstepcounter{questao}\medskip\noindent
  \textbf{Question \thequestao.}~#1\par}
\newenvironment{alternativas}{\begin{enumerate}[label=\Alph*.,leftmargin=2.2em,topsep=1pt]}{\end{enumerate}}
\newcommand{\gabarito}[1]{\ifsolucoes\par\noindent\textbf{Answer:}~#1\fi}
\newcommand{\resolespaco}{\ifsolucoes\else\par\textit{Solution:}\vspace{\espacoresol}\par\fi}
\newcommand{\passo}[2]{\par\smallskip\noindent\textbf{Step --- #1.}~#2}
\newcommand{\intuicao}[1]{\par\smallskip\noindent\textit{Economic intuition.}~#1}
\newcommand{\refbib}[1]{\par\smallskip\noindent{\footnotesize\color{black!60}\textbf{Ref:} #1}}
\newcommand{\rubrica}[1]{\par\smallskip\noindent\textbf{Grading rubric.}\begin{itemize}[leftmargin=1.4em,topsep=1pt]#1\end{itemize}}

\newcommand{\cabecalho}{%
\begin{center}
  {\Large\bfseries \disciplina}\\[2pt]
  {\large Practice list: \listatema}\ifsolucoes\\[2pt]\textit{Solutions and grading rubric}\fi
\end{center}
\vspace{2pt}
{\footnotesize\noindent\textbf{Instructions.} Multiple choice: mark the single correct option (no
penalty for wrong answers). \emph{True/False}: every 2 wrong items cancel 1 right one
(score $=\max\{0,\text{right}-\text{wrong}/2\}\times$ item value); blank items do not count.
Discursive questions: show every step. Numbers are illustrative and chosen to give clean answers.\par}}
.
\usepackage[utf8]{inputenc}
\usepackage[T1]{fontenc}
\usepackage{lmodern}   % Type 1 vector fonts (replaces PK bitmaps)
\usepackage{cmap}      % ToUnicode CMap -> copyable text and accents
\usepackage[english]{babel}
\usepackage{amsmath,amssymb,amsfonts}
\usepackage[a4paper,margin=2.2cm]{geometry}
\usepackage{enumitem}
\usepackage{fancyhdr}
\usepackage{comment}
\usepackage{xcolor}
\usepackage{microtype}
\usepackage{booktabs}
\usepackage{tikz}

\newif\ifsolucoes
\newcommand{\disciplina}{Macroeconomics I --- MPE Insper}
\newcommand{\espacoresol}{4.5cm}   % height of the answer space in the blank list
% The sol environment must be decided after the list sets its toggle.
\AtBeginDocument{\ifsolucoes\includecomment{sol}\else\excludecomment{sol}\fi}

\pagestyle{fancy}\fancyhf{}
\lhead{\footnotesize\disciplina}
\rhead{\footnotesize Practice list\ifsolucoes{} --- Solutions and rubric\fi}
\cfoot{\footnotesize\thepage}
\renewcommand{\headrulewidth}{0.3pt}
\setlength{\headheight}{14pt}
\setlist{topsep=2pt,itemsep=2pt}
\newcounter{questao}

\newcommand{\bloco}[3]{\bigskip\noindent
  {\large\bfseries\color{blue!55!black}Block #1 --- #2}\hfill\textit{(#3)}\par
  \rule{\linewidth}{0.3pt}\smallskip}
\newcommand{\questao}[1]{\refstepcounter{questao}\medskip\noindent
  \textbf{Question \thequestao.}~#1\par}
\newenvironment{alternativas}{\begin{enumerate}[label=\Alph*.,leftmargin=2.2em,topsep=1pt]}{\end{enumerate}}
\newcommand{\gabarito}[1]{\ifsolucoes\par\noindent\textbf{Answer:}~#1\fi}
\newcommand{\resolespaco}{\ifsolucoes\else\par\textit{Solution:}\vspace{\espacoresol}\par\fi}
\newcommand{\passo}[2]{\par\smallskip\noindent\textbf{Step --- #1.}~#2}
\newcommand{\intuicao}[1]{\par\smallskip\noindent\textit{Economic intuition.}~#1}
\newcommand{\refbib}[1]{\par\smallskip\noindent{\footnotesize\color{black!60}\textbf{Ref:} #1}}
\newcommand{\rubrica}[1]{\par\smallskip\noindent\textbf{Grading rubric.}\begin{itemize}[leftmargin=1.4em,topsep=1pt]#1\end{itemize}}

\newcommand{\cabecalho}{%
\begin{center}
  {\Large\bfseries \disciplina}\\[2pt]
  {\large Practice list: \listatema}\ifsolucoes\\[2pt]\textit{Solutions and grading rubric}\fi
\end{center}
\vspace{2pt}
{\footnotesize\noindent\textbf{Instructions.} Multiple choice: mark the single correct option (no
penalty for wrong answers). \emph{True/False}: every 2 wrong items cancel 1 right one
(score $=\max\{0,\text{right}-\text{wrong}/2\}\times$ item value); blank items do not count.
Discursive questions: show every step. Numbers are illustrative and chosen to give clean answers.\par}}

\solucoesfalse                 % <<<< \solucoestrue for the ANSWER KEY
\newcommand{\listatema}{The New Keynesian AS--AD Model (Lista 7)}

\begin{document}
\cabecalho

{\footnotesize\noindent\textbf{Notation (Benigno 2015).} Lower-case letters are logs. AD:
$y=\bar y_n+(g-\bar g)-\sigma[\,i-(\bar p-p)-\rho\,]$ (eq.\ 21, no consumption taxes).
AS: $p-p^e=\kappa(y-y_n)$ (eqs.\ 17, 20). Natural output
$y_n=\frac{1+\eta}{\sigma^{-1}+\eta}a+\frac{\sigma^{-1}}{\sigma^{-1}+\eta}g-\frac{1}{\sigma^{-1}+\eta}\mu$ (eq.\ 15);
efficient output $y_e$ is the same without the $\mu$ term (eq.\ 19). Bars denote the long run.\par}

% =============================================================================
\bloco{1}{Foundations: AD, AS, natural and efficient output}{Benigno \S\S1--5 \,\textbullet\, 25 points}

\questao{(10 points) Objective items.}

\textbf{a)} (3 points) If the share of firms whose prices are fixed in advance rises, then:
\begin{alternativas}
\item $\kappa$ rises, the AS gets steeper, and a cut in $i$ moves output by more.
\item $\kappa$ falls, the AS gets flatter, and a cut in $i$ moves output by more.
\item $\kappa$ is unchanged, the AS shifts up, and a cut in $i$ moves prices by more.
\item $\kappa$ falls, the AS gets steeper, and a cut in $i$ moves prices by more.
\end{alternativas}
\gabarito{B}
\begin{sol}
\passo{slope}{Only flexible-price firms respond to marginal cost; with fewer of them the price
level responds less to the output gap: $\kappa\downarrow$, a flatter AS, so shifts in AD
move $y$ more and $p$ less.}
\refbib{Benigno (2015), \S4 and \S4.1, pp.~506--507.}
\end{sol}

\medskip
\textbf{b)} (3 points) With $\sigma=\eta=1$, a temporary rise of $0.04$ in the (log) desired markup
$\mu$ changes $(y_n,\,y_e)$ by:
\begin{alternativas}
\item $(-0.04,\,-0.04)$
\item $(0,\,-0.02)$
\item $(-0.02,\,0)$
\item $(-0.02,\,-0.02)$
\end{alternativas}
\gabarito{C}
\begin{sol}
\passo{coefficients}{$\partial y_n/\partial\mu=-1/(\sigma^{-1}+\eta)=-\tfrac12$, so $\Delta y_n=-0.02$;
$y_e$ does not depend on $\mu$.}
\refbib{Benigno (2015), \S4.4, eqs.\ (15) and (19), p.~508.}
\end{sol}

\medskip
\textbf{c)} (2 points, T or F) In Benigno's model the AD slopes down because a higher price level
raises money demand and therefore the nominal interest rate.
\gabarito{F}
\begin{sol}
\passo{channel}{There is no LM curve: the central bank sets $i$. Given $\bar p$, a higher $p$ lowers
expected inflation $\bar p-p$, raises the real rate $i-(\bar p-p)$, and households save more:
an intertemporal channel through the Euler equation.}
\refbib{Benigno (2015), \S3, eq.\ (21).}
\end{sol}

\medskip
\textbf{d)} (2 points, T or F) After a permanent productivity increase, with $i$ unchanged, the
economy ends up at $p=p^e$ and $y=y_n$: both curves shift and no policy action is needed.
\gabarito{T}
\begin{sol}
\passo{\S6.2}{$y_n$ and $\bar y_n$ rise by the same amount: AS moves right through $(p^e,y_n')$ and
AD moves right by $\Delta\bar y_n$ (consumption smoothing). The natural rate is unchanged.}
\refbib{Benigno (2015), \S6.2, p.~512.}
\end{sol}

% -----------------------------------------------------------------------------
\questao{(15 points) Discursive. Copom's model. Take $\sigma=\eta=1$, $\kappa=0.5$, $\rho=0.02$,
$p^e=\bar p=0$, $\bar y_n=0$, $g=\bar g$, and initially $a=\mu=0$, so $i=\rho$ gives $y=p=0$.}

\textbf{a)} (5 points) Solve AD and AS for equilibrium $y$ and $p$ as functions of $i$, $\bar p$,
$p^e$, $y_n$ and $\bar y_n$. Which variables move AD and which move AS?
\resolespaco
\begin{sol}
\passo{AD}{$y=X-\sigma p$ with $X\equiv\bar y_n+(g-\bar g)-\sigma(i-\rho)+\sigma\bar p$.}
\passo{solve}{Substitute $p=p^e+\kappa(y-y_n)$: $y=\dfrac{X-\sigma p^e+\sigma\kappa\,y_n}{1+\sigma\kappa}$,
then $p=p^e+\kappa(y-y_n)$.}
\passo{shifters}{AD: $i$, $g$, $\bar g$, $\bar p$, $\bar y_n$ (future $\bar a$, $\bar\mu$). AS: $p^e$ and
$y_n$ (current $a$, $g$, $\mu$).}
\refbib{Benigno (2015), \S4--5.}
\rubrica{
\item Rewrites AD with $X$: 1 pt.
\item Solves for $y$ and $p$: 2.5 pts.
\item Correct list of shifters: 1.5 pts. Putting $\bar p$ on the AS: $-1$ pt.
}
\end{sol}

\textbf{b)} (5 points) A temporary productivity gain: $a=0.03$, $\bar a=0$. Compute $y_n$, then $y$,
$p$ and the gap with $i=\rho$. Which $i$ restores $p=p^e$ and $y=y_n$? Show that it violates the zero
lower bound, compute the outcome at $i=0$, and find the long-run price level $\bar p$ that the
central bank must credibly promise to reach the target with $i=0$.
\resolespaco
\begin{sol}
\passo{natural}{$y_n=\frac{2}{2}\cdot0.03=0.03$; $y_e=y_n$.}
\passo{no action}{$X=0$: $y=\frac{0.5\cdot0.03}{1.5}=0.01$, $p=0.5(0.01-0.03)=-0.01$; gap $=-0.02$.}
\passo{required $i$}{Need $y=0.03$, $p=0$, so $X=0.03$: $-(i-0.02)=0.03\Rightarrow i=-1\%<0$.}
\passo{at the ZLB}{$X=0.02$: $y=\frac{0.02+0.015}{1.5}=0.0233$, $p=-0.0033$.}
\passo{expectations}{With $i=0$ we need $\sigma\bar p=0.01$: $\bar p=1\%$, giving $y=0.03$, $p=0$.}
\intuicao{The economy can produce more but demand does not rise with it today; the real rate must
fall. When $i$ cannot, a higher promised future price level raises expected inflation and
lowers the real rate, but only if the promise is believed.}
\refbib{Benigno (2015), \S6.1 (p.~512) and \S9.}
\rubrica{
\item $y_n$ and outcome with $i=\rho$: 1.5 pts. Shifting AD in a \emph{temporary} shock: max.\ 2 pts.
\item Required $i$ and the ZLB: 1.5 pts.
\item Outcome at $i=0$ and the required $\bar p$: 2 pts.
}
\end{sol}

\textbf{c)} (5 points) A temporary markup shock (oil, a tax on firms): $\mu=0.04$, $a=0$. Compute
$y_n$, $y_e$ and the outcome with $i=\rho$. Find $i$ under strict price stability. The bank minimizes
$\mathcal L=(p-p^e)^2+\lambda(y-y_e)^2$ subject to the AS, with $\lambda=0.25$: derive the optimal
$y$, $p$ and $i$. Draw the diagram.
\resolespaco
\begin{sol}
\passo{levels}{$y_n=-0.02$, $y_e=0$. With $i=\rho$: $y=\frac{0.5(-0.02)}{1.5}=-0.0067$,
$p=0.5(-0.0067+0.02)=+0.0067$: stagflation.}
\passo{strict}{$p=0\Rightarrow y=y_n=-0.02$, so $X=-0.02$ and $i=\rho+0.02=4\%$.}
\passo{optimal}{Substitute the AS: $\min_y\kappa^2(y-y_n)^2+\lambda y^2\Rightarrow
y=\dfrac{\kappa^2}{\kappa^2+\lambda}y_n=\tfrac12(-0.02)=-0.01$, $p=0.5(-0.01+0.02)=0.005$.
From (a): $-0.015=X-0.01\Rightarrow X=-0.005\Rightarrow i=2.5\%$.}
\passo{trade-off}{Unlike (b), $y_n\ne y_e$: the bank cannot close the efficient gap and keep
prices stable at once. It splits the shock according to $\lambda$.}
\begin{center}
\begin{tikzpicture}[scale=1.7,>=stealth,font=\small]
  \draw[->] (-3,0) -- (2,0) node[right]{$y$ ($\times100$)};
  \draw[->] (0,-1.5) -- (0,2) node[above]{$p$ ($\times100$)};
  \draw[thick] (-1.6,1.6) -- (1.4,-1.4) node[below right]{AD ($i=\rho$)};
  \draw[thick,dashed] (-2.4,1.9) node[above]{AD$''$ ($i=2.5\%$)} -- (0.8,-1.3);
  \draw[thick,blue] (-3,-1.5) -- (2,1) node[right]{AS};
  \draw[thick,red] (-3,-0.5) -- (1.2,1.6) node[right]{AS$'$};
  \draw[dotted] (-2,0) -- (-2,-1.2) node[below]{$y_n'$};
  \fill (0,0) circle (1.2pt) node[below right]{$E$};
  \fill (-0.667,0.667) circle (1.2pt) node[above right]{$E'$};
  \fill (-1,0.5) circle (1.2pt) node[below right]{$E''$};
\end{tikzpicture}\\[2pt]
{\footnotesize $E=(y_e,p^e)$ is the origin; $E'$: no policy; $E''$: optimal policy; strict price
stability would sit at $(y_n',p^e)$.}
\end{center}
\refbib{Benigno (2015), \S7 (p.~513) and \S11.}
\rubrica{
\item $y_n$, $y_e$ and outcome at $i=\rho$: 1.5 pts.
\item Strict price stability: 1 pt.
\item Optimal $y$, $p$, $i$: 1.5 pts. Measuring the gap against $y_n$ in the loss: max.\ 2 pts in the item.
\item Diagram with axes, AS shift through $(p^e,y_n')$, AD shift and points: 1 pt.
}
\end{sol}

% =============================================================================
\bloco{2}{Policy analysis}{Benigno \S\S6--12 \,\textbullet\, 25 points}

\questao{(10 points) Objective items.}

\textbf{a)} (3 points) The Focus survey suddenly expects higher \emph{future} productivity; current
productivity is unchanged (Benigno \S6.3). The policy that keeps $p=p^e$ and $y=y_n$ is to:
\begin{alternativas}
\item raise $i$, since AD shifts right while AS stays put.
\item cut $i$, since AS shifts right while AD stays put.
\item keep $i$, since $y_n$ today is unchanged by news.
\item promise a higher $\bar p$, since AD shifts left today.
\end{alternativas}
\gabarito{A}
\begin{sol}
\passo{\S6.3}{Higher $\bar a$ raises $\bar y_n$: households want to consume more now (AD right).
Current $y_n$ is unchanged (AS fixed), so output rises above both $y_n$ and $y_e$: the bank must
\emph{tighten}.}
\refbib{Benigno (2015), \S6.3, p.~513.}
\end{sol}

\medskip
\textbf{b)} (3 points) With $\sigma=\eta=1$, $\kappa=0.5$ and $i$ fixed, a temporary rise of $0.01$ in $g$
raises output by:
\begin{alternativas}
\item $0.0100$
\item $0.0050$
\item $0.0150$
\item $0.0083$
\end{alternativas}
\gabarito{D}
\begin{sol}
\passo{both curves move}{AD: $X$ rises by $0.01$. AS: $y_n$ rises by $\frac{\sigma^{-1}}{\sigma^{-1}+\eta}0.01=0.005$.
$y=\frac{0.01+0.5\cdot0.005}{1.5}=0.0083$: multiplier $0.83<1$; the gap rises only $0.0033$.}
\refbib{Benigno (2015), \S8.}
\end{sol}

\medskip
\textbf{c)} (2 points, T or F) At the zero lower bound, monetary policy has no instrument left to
raise current output.
\gabarito{F}
\begin{sol}
\passo{expectations}{A credible commitment to a higher long-run price level $\bar p$ raises expected
inflation, lowers the real rate with $i=0$ and shifts AD right.}
\refbib{Benigno (2015), \S9.}
\end{sol}

\medskip
\textbf{d)} (2 points, T or F) Productivity and government-spending shocks move $y_n$ and $y_e$ by
the same amount, so they create no trade-off between stabilizing prices and the efficient gap.
\gabarito{T}
\begin{sol}
\passo{eqs.\ (15), (19)}{The coefficients on $a$ and $g$ are identical; only $\mu$ separates $y_n$
from $y_e$. Under flexible wages, ``divine coincidence'' holds for these shocks.}
\refbib{Benigno (2015), \S4.4 and \S11.}
\end{sol}

% -----------------------------------------------------------------------------
\questao{(15 points) Discursive.}

\textbf{a)} (5 points) Derive $y_n$ and $y_e$ from first-order conditions in log form. Use the
labor supply condition $\eta l+\tilde\sigma^{-1}c=w-p$, the price rule of flexible firms
$p=w-a+\mu$, the technology $y=a+l$ and the resource constraint $y=s_cc+g$, with
$\sigma=\tilde\sigma s_c$. For $y_e$, replace the price rule by the planner's condition
``marginal rate of substitution $=$ marginal product of labor''. Explain why a markup shock
moves $y_n$ but not $y_e$.
\resolespaco
\begin{sol}
\passo{substitute}{$c=(y-g)/s_c\Rightarrow\tilde\sigma^{-1}c=\sigma^{-1}(y-g)$; $l=y-a$.}
\passo{natural}{Real wage $w-p=a-\mu$: $\eta(y-a)+\sigma^{-1}(y-g)=a-\mu$, so
$y_n=\frac{(1+\eta)a+\sigma^{-1}g-\mu}{\sigma^{-1}+\eta}$.}
\passo{efficient}{Planner: MRS $=$ MPL (log MPL $=a$): $\eta(y-a)+\sigma^{-1}(y-g)=a$, the same
without $\mu$.}
\passo{why}{The markup is a wedge between the real wage and the marginal product of labor; the
planner faces no such wedge. Monopoly power and distorting taxes lower natural output without
changing what is efficient: $y_n-y_e=-\mu/(\sigma^{-1}+\eta)$.}
\refbib{Benigno (2015), \S\S3--4.4, eqs.\ (15) and (19).}
\rubrica{
\item Substitutions: 1 pt.
\item Derives $y_n$: 1.5 pts.
\item Derives $y_e$: 1.5 pts.
\item Wedge explanation: 1 pt. ``The markup is inefficient'' with no wedge argument: 0.5 pt.
}
\end{sol}

\textbf{b)} (5 points) The Copom cuts the Selic by one percentage point. Compare the short- and
long-run effects predicted by (i) the New Classical model, (ii) the original Keynesian (IS--LM)
model and (iii) Benigno's New Keynesian model, along three dimensions: price adjustment and
expectations; the policy instrument and how AD is built; and whether the model can rank outcomes
by welfare.
\resolespaco
\begin{sol}
\passo{New Classical}{Flexible prices, rational expectations, competitive markets. An anticipated
cut is neutral; only surprises move output (Lucas). No welfare-relevant inefficiency to correct.}
\passo{Keynesian}{Fixed or adaptive prices, AD from IS--LM with the money supply as the
instrument. The cut raises output persistently; the model is not built from optimizing agents and
cannot rank outcomes by household welfare.}
\passo{New Keynesian}{Rational expectations, monopolistic competition, a fraction of prices preset.
The instrument is $i$; AD comes from the Euler equation. Short run: AD shifts right, $y$ and $p$
rise along the AS. Long run: all prices adjust, $y=y_n$, money is neutral. Welfare is the
household's utility, approximated by the quadratic loss in $(p-p^e)$ and $(y-y_e)$.}
\refbib{Benigno (2015), \S1 and \S12; Lista 7, Q1; Carlin \& Soskice (2024) for the IS--LM contrast.}
\rubrica{
\item Each framework with the three dimensions: 1.5 pts each (4.5 pts).
\item Short- vs long-run distinction in the NK model: 0.5 pt.
\item Building NK AD from LM / money demand: max.\ 3 pts in the item.
}
\end{sol}

\textbf{c)} (5 points) A deleveraging shock (Eggertsson--Krugman): the debt limit of borrowers
falls sharply. Explain why the economy can fall into a liquidity trap, why the fiscal multiplier
becomes large, and why more price flexibility can make the recession worse (paradox of
flexibility). Under what condition on long-run prices does the paradox disappear?
\resolespaco
\begin{sol}
\passo{trap}{Borrowers must cut consumption; AD falls so much that the natural real rate turns
negative. With $i\ge0$ the bank cannot lower the real rate enough: output stays below $y_n$.}
\passo{multiplier}{Ricardian equivalence fails: borrowers are constrained and spend extra income,
and at $i=0$ there is no crowding-out through higher rates. The multiplier can exceed one
(Benigno reports 1.29 to 1.62), against the upper bound of one in \S8.}
\passo{flexibility}{Debt is nominal. If inflation between the short and the long run is fixed (as in
Eggertsson--Krugman), a lower $p$ also lowers $\bar p$ and raises the real value of debt: borrowers
cut consumption, and the AD can slope \emph{up}. A steeper AS (more flexible prices) then means a
larger fall in prices and a larger fall in output. For the same reason a downward AS shift (higher
productivity, lower markups or taxes) becomes contractionary: the paradox of toil.}
\passo{caveat}{If the long-run price level is anchored whatever happens to current prices, the AD
slopes down again and both paradoxes disappear (Benigno, footnote 20).}
\intuicao{In a debt deflation, falling prices transfer resources to savers with a low propensity to
spend, away from borrowers who must cut back.}
\refbib{Benigno (2015), \S9 and \S10, eq.\ (32), pp.~518--521.}
\rubrica{
\item Liquidity trap mechanism with negative natural rate: 2 pts.
\item Large multiplier with both reasons (constrained borrowers, no crowding-out): 1.5 pts.
\item Paradox of flexibility through the real value of debt: 1.5 pts. Arguing through the
  real interest rate with a fixed $\bar p$: max.\ 0.5 pt on this line.
}
\end{sol}

\end{document}
